\documentclass[11pt,a4paper]{article}
\usepackage{gtu-el-yazisi}
\begin{document}
\notbilgi{RF Tasarım}{RF Link, Antenna Gain, Path Loss, Link Margin}{Ders9.pdf}

\sayfa{1}
\section{RF link}

\elcizimi[0.7]{sekiller/Ders9-s1-1.png}[1]

\textbf{1) Antenna Gain}

Antenna efficiency \quad $\eta = \dfrac{P_T}{P_{\mathrm{in}}} < 1$

\[ G_A = \frac{\text{Maximum power density of the actual antenna}}{\text{power density of isotropic antenna}} \]

\elcizimi[0.6]{sekiller/Ders9-s1-2.png}[1]

\elcizimi[0.35]{sekiller/Ders9-s1-3.png}[1]

eğer $P_T$ gücünü yayıyorsa

\[ \underbrace{P_D}_{\text{power density}} = \frac{P_T}{4\pi d^2} \]

Antenler her yöne eşit dağılım yapmaz. Bu nedenle her yöne eşit dağılım yaptığını varsaydığımız isotropic antenle mukayese ederiz

\[ G_A \;\rightarrow\; P_D = \frac{P_T\cdot \underbrace{G_A}_{\substack{\text{if antenna}\\ \text{has gain}}}}{4\pi d^2} \]

if antenna has Gain and efficiency

\[ \rightarrow\; P_D = \frac{P_T\,G_A\,\eta}{4\pi d^2}\ [\mathrm{W/m^2}] \]

$\rightarrow$ Güç yoğunluğu mesafenin karesiyle ters orantılı olarak değişiyor.

\subsection{Alıcıdaki:}

for

\[ R_x \;\rightarrow\; A_R = \frac{G_R\cdot\overset{\text{wave length}}{\lambda}{}^{2}}{4\pi}\ [\mathrm{m^2}] \;\rightarrow\; \underset{\text{(Etkin açıklık)}}{\text{Effective Arpeture Size}} \]

\[ P_R = P_D\,A_r \]

\[ P_R = \frac{P_T\,G_A\cdot\eta_A}{4\pi d^2}\cdot\frac{G_R\cdot\lambda^2\,(\eta_r)}{4\pi} \]

\[ P_R = P_T\cdot G_A\cdot G_R\cdot\eta_A\cdot\eta_R\left(\frac{\lambda}{4\pi d}\right)^2 \]

Tx antenna: $G_A$, $\eta_A$ \qquad Rx antenna: $G_R$, $\eta_R$

\sayfa{2}
\begin{align*}
P_R\big|_{\mathrm{dBm}} &= P_T\big|_{\mathrm{dBm}} \\
&\quad + \underbrace{G_A\big|_{\mathrm{dB}} + G_R\big|_{\mathrm{dB}} + \eta_A\big|_{\mathrm{dB}} + \eta_R\big|_{\mathrm{dB}} + \overbrace{20\log\left(\frac{\lambda}{4\pi d}\right)}^{L_{\mathrm{path,loss(LOS)}}}}_{\substack{L_{\mathrm{link\,dB}}(LOS)\\ \text{Line of side}}}
\end{align*}

$\rightarrow$ Alıcıyla vericinin birbirini gördüğü durum için

\begin{ornek}
point-point (LOS) \\
$G_R = 60$ dB \\
$f = 60$ GHz \\
$P_{\mathrm{density}} = 1\ \mathrm{pW/cm^2}$

what is the power at output of antenna connected to RF electronics

\elcizimi[0.35]{sekiller/Ders9-s2-1.png}[2]
\end{ornek}

\[ \lambda = \frac{c}{f} \]

\[ A_R = \frac{G_R\cdot\lambda^2}{4\pi} = \frac{10^6\left(\dfrac{3\cdot 10^8}{60\cdot 10^9}\right)^2}{4\pi} = 1{,}99\ \mathrm{m^2} \approx 2\ \mathrm{m^2} \]

\[ P_R = P_D\cdot A_R = \frac{10^{-12}}{\mathrm{cm^2}}\cdot 2\cdot 10^4\ \mathrm{cm^2} = 20\ \mathrm{nW} \]

\section{LOS olmazsa}

\elcizimi[0.5]{sekiller/Ders9-s2-2.png}[2]

direct signal is dominant

Ricean fading $\left[\dfrac{1}{d}\right]$

\elcizimi[0.45]{sekiller/Ders9-s2-3.png}[2]

\[ \rightarrow\ \frac{\text{reflected waves}}{\text{Rayleigh fading}}\ \left[\frac{1}{d^4}\right] \]

\[ L_{\mathrm{path,loss}}(LOS) = \left(\frac{4\pi d}{\lambda}\right)^2 \]

\[ L_{\mathrm{path,loss}}\underbrace{(NLOS)}_{\substack{\text{non line}\\ \text{of side}}} = \left(\frac{4\pi d}{\lambda}\right)^2 (d)^{n-2} \]

$n$: number of paths

\[ L_{\mathrm{path,NLOS}}\big|_{\mathrm{dB}} = 20\log\left(\frac{4\pi d}{\lambda}\right) + 10(n-2)\log(d) \]

\[ = 10n\log(d) + 20\log\left(\frac{4\pi}{\lambda}\right) \]

\sayfa{3}
\begin{ornek}
$f = 5{,}6$ GHz \\
$G_A = 35$ dB \\
$G_R = 6$ dB \\
$d = 200$ m \\
$\dfrac{1}{d^3} \rightarrow n = 3$
\end{ornek}

\begin{align*}
L_{\mathrm{Link}} &= -G_A - G_R + 10n\log(d) + 20\log\left(\frac{4\pi}{\lambda}\right) \\
&= -35 - 6 + 10\cdot 3\log(100)\footnotemark + 20\log\left(\frac{4\pi}{\dfrac{3\cdot 10^8}{5{,}6\cdot 10^9}}\right) \\
&= 75{,}4\ \mathrm{dBm}
\end{align*}
\footnotetext{Claude notu: $d = 200$ m olduğundan muhtemelen $\log(200)$ olmalı (75,4 sonucu buna uyuyor).}

\section{Link Margin}

\[ P_R = P_T - L_{\mathrm{Link}} \]

\[ \underbrace{LM}_{\text{Link Margin}} = P_R - \overbrace{S}^{\text{alıcı hassasiyeti}} \]

\begin{align*}
S &= -100\ \mathrm{dBm} \\
P &= -97\ \mathrm{dBm} \\
LM &= -97 - (-100) \\
LM &= 3\ \mathrm{dB}
\end{align*}

Bu linkte oluşabilecek bozulmalara 3 dBm bir pay bulunmaktadır. Eğer 5 dBm'lik bir bozulma olursa alıcı artık çalışmıyordur ($3-5 = -2$ dBm)

\section{Link Budget Table}

\elcizimi[0.5]{sekiller/Ders9-s3-1.png}[3]

\end{document}
